\emph{Trace-channel response of the fibre part: the \(\vartheta\) response of the first term (conditional).} To the hypotheses of the set of preceding paragraphs (equal weights, \(V=T\), frozen symbol, product background, one loop) add the pure-trace channel convention of the preceding paragraph, \(\bar h_{\mu\nu}=(t/4)\bar g_{\mu\nu}\), \(t=h^\mu{}_\mu\). Of the data entering the level form of the fibre flux, the wave in this channel moves only the scalar \(\vartheta=\coth^2\tau=-y\cdot y/m^2\): the contraction is taken in the fluctuated metric \(\bar g+\bar h\) with the fibre coordinate and \(m\) held, so \(\vartheta\to(1+t/4)\vartheta\) exactly, while \(\mu_J\), \(c_J=J(J+1)\), \(M_A^2=\chi^2\mu_A^2\) and \(\mathcal D=M_A^2+\mu_J\) do not move. In this channel \(\bar h\) is not negligible along \(y\), so the recorded rigidity of the densities of the class in the static channel at large \(\tau\) does not reach it. The level factor in the fibre-flux display is a sum of terms of degrees \(2,1,0\) in \(\vartheta\) with coefficients depending on \(\mathcal D\) and \(c_J\) alone: write that bracket as \(B_0\); the first term responds as \((1+t/4)^2B((1+t/4)\vartheta,\mathcal D,c_J)\), its coefficient of \(t\) at \(q=0\) is per level \(B_0/2+\vartheta\,\partial_\vartheta B/4\), the first part the determinant profile of the preceding paragraph and the extra tadpole \(\vartheta(15\vartheta\mathcal D-2c_J)/(8\mathcal D^2)\), and its coefficient of \(t^2\) is \(B_0/16+(4\vartheta\,\partial_\vartheta B+\vartheta^2\partial_\vartheta^2B)/32\), the first part that of \(\sqrt{|g|}\) alone and the extra quadratic part \(\vartheta(75\vartheta\mathcal D-8c_J)/(64\mathcal D^2)\). Applying the proper-time measure \(\mathrm ds/s\) to the recorded integrand average \(s^2[\tfrac{15}4\vartheta^2-s\vartheta c_J+s^2c_J^2]\) instead gives the summand \(\tfrac{15\vartheta^2}{4\mathcal D^2}-\tfrac{2\vartheta c_J}{\mathcal D^3}+\tfrac{6c_J^2}{\mathcal D^4}\), whose sign is that of a positive definite quadratic form in \(\vartheta\mathcal D\) and \(c_J\) and whose extra parts are \(\vartheta(15\vartheta\mathcal D-4c_J)/(8\mathcal D^3)\) and \(\vartheta(75\vartheta\mathcal D-16c_J)/(64\mathcal D^3)\): same structure, same sign conditions. The \(\vartheta\)-linear part of \(U_2\) in the potential display responds likewise as \((1+t/4)^3\) against the profile \((1+t/4)^2\). Within these hypotheses, then, the response of the fibre levels to \(\bar h\) in \(v^M-2\sigma^M\) does not vanish in the pure-trace \(q=0\) channel: it is carried by \(\vartheta\), and it supplies both an uncancelled tadpole and a quadratic part beyond the determinant profile, both positive at each level when \(15\vartheta\mathcal D>4c_J\), a condition that also implies both quadratic conditions \(75\vartheta\mathcal D>16c_J\) and \(75\vartheta\mathcal D>8c_J\). The candidate flagged above as able to supply \(\mathcal M\) without cancelling its own tadpole is therefore present at the level form, and the rigidity recorded above is confined to the terms carrying \(\sqrt{|g|}\) alone, of which the first term with its fibre flux is not one. Not settled here: the sums over \(A\) and \(J\le n\) at the stated cut-off \(\mathcal D\gg\vartheta\), the response of the flux insertions and of the weights to the wave, the corrections of relative order \(s\) and \(s\vartheta\) recorded there, the second-order non-derivative terms at \(\chi^2\lesssim\vartheta\), and the comparison of magnitudes with \(q_b^2\), which remains convention-bound as recorded, so these identities do not fix the sign of the full \(\mathcal M\) and do not decide between the Yukawa threshold and the tachyonic pole recorded above. These statements are for the hypotheses and the channel convention stated here; they retire the recorded item on the response of the fibre levels to \(\bar h\) only for this channel component, they change no status label, and Problems~C, F and~N remain open and PROGRAMME-CONDITIONAL.
